\paragraph{C.3o (Proved/refuted): the remaining ordinary degree-one boundary dimensions over $\mathbb C$.}
In the coefficient-field setting of C.3n, with ordinary
algebraic equivalence and rational coefficients, the
universal injectivity threshold over $\mathbb C$ is
$N\ge3$. Surjectivity is automatic for $N\le2$ and
fails in every dimension $N\ge3$. The assertions
for $N\ge4$ follow from C.3m--n. Here are explicit
examples and proofs for the remaining dimensions;
they use actual line bundles on strictly henselian
punctures.

\begin{center}
\begin{tabular}{lll}
\toprule
Dimension $N$ & Always injective & Always surjective \\
\midrule
$1$ & Yes, both groups vanish & Yes, both groups vanish \\
$2$ & No, the cone example below & Yes, the target vanishes \\
$\ge3$ & Yes, by C.3n & No, below and C.3m \\
\bottomrule
\end{tabular}
\end{center}

\begin{proof}
Let
\[
 F=z_0^4+z_1^4+z_2^4,\qquad
 C=V_+(F)\subset\mathbb P^2_{\mathbb C},\qquad
 L=\mathcal O_C(1).
\]
This is the smooth plane quartic already used in
C.3i and C.3m, with genus three and $\deg L=4$.
Define the two-dimensional excellent normal
strictly henselian local ring
\[
 R_2=
 \left(\mathbb C[z_0,z_1,z_2]/(F)\right)_{(z_0,z_1,z_2)}^{\mathrm{sh}},
 \qquad
 X_2=\operatorname{Spec}R_2\setminus\{\mathfrak m_2\}.
\]
The affine cone is smooth away from its vertex,
so $X_2$ is regular. Normality follows, for example,
from the hypersurface Cohen--Macaulay property and
regularity in codimension one; excellence preserves
normality under strict henselization. As throughout,
the algebraically closed residue field is identified
with the coefficient field $\mathbb C$.

The cone projection gives an actual morphism
$\rho:X_2\to C$. We claim that
\begin{equation}\label{eq:C3-low-dimension-Jacobian}
 \rho^*:\operatorname{Pic}^0(C)(\mathbb C)
       \lhook\joinrel\longrightarrow
       \operatorname{Pic}(X_2)_{\mathrm{alg}}
\end{equation}
is injective. Pulling back the usual Poincare
families on $C$ shows that its image consists of
ordinary algebraically trivial classes.
For injectivity, base change the proper cone
resolution $\operatorname{Tot}(L^{-1})\to
\operatorname{Spec}\mathbb C[z_0,z_1,z_2]/(F)$
to $\operatorname{Spec}R_2$. The resulting regular
scheme $Z_2$ has exceptional divisor $C$, is
isomorphic to $X_2$ off $C$, and satisfies
\[
                 \mathcal O_{Z_2}(C)|_C=L^{-1}.
\]
If $Q\in\operatorname{Pic}^0(C)$ pulls back to
zero on $X_2$, its pullback to $Z_2$ has a rational
trivialization with divisor supported on the single
prime divisor $C$. Thus it is $\mathcal O_{Z_2}(bC)$
for an integer $b$. Restriction to $C$ gives
$Q\simeq L^{-b}$. Degrees give $0=-4b$, hence
$b=0$ and $Q=0$. This proves
\eqref{eq:C3-low-dimension-Jacobian} on the strict
henselization itself.

An explicit infinite-order class makes the
rational-coefficient failure just as definite.
Choose $u\in\mathbb C$ transcendental over
$\overline{\mathbb Q}$, choose $v$ with
$v^4=-1-u^4$, put $p=(u:v:1)$, and fix
$c_0=(1:\zeta:0)$ with $\zeta^4=-1$. Then
\[
                         Q=\mathcal O_C(p-c_0)
\]
has infinite order. Indeed $C$, $c_0$, its
Jacobian, and the Abel closed immersion are
defined over $\overline{\mathbb Q}$. Every
torsion point of that Jacobian over $\mathbb C$
is already defined over $\overline{\mathbb Q}$,
since each torsion subgroup scheme is finite
over this algebraically closed field. A torsion
$Q$ would therefore force $p$ to be defined
over $\overline{\mathbb Q}$, contrary to the
choice of $u$. The Abel closed immersion and
the finite torsion property are the standard
Jacobian facts recorded in
\citet[Propositions 2.1--2.3]{D_MilneJacobians2021}
and \citet[Section 12]{D_MilneAbelian2022}.

\emph{Failure of injectivity for $N=2$.}
Take $f=z_2\in R_2$. The divisor puncture
\[
                         Y_2=X_2\cap V(z_2)
\]
is regular and zero dimensional: before
henselization its four branches are the
distinct lines $z_0^4+z_1^4=0$, meeting only
at the removed vertex. Equivalently, the
finite normalization of $R_2/(z_2)$ is a
product of four strictly henselian discrete
valuation rings, and their punctures are
four spectra of fields. Hence
$\operatorname{Pic}(Y_2)=0$. The infinite-order
class $\rho^*Q$ gives a nonzero element of
$J^1(X_2)$ restricting to zero in $J^1(Y_2)$.

\emph{Failure of surjectivity for $N=3$.}
Now take
\[
 R_3=\mathbb C[z_0,z_1,z_2]_{(z_0,z_1,z_2)}^{\mathrm{sh}},
 \qquad X_3=\operatorname{Spec}R_3\setminus\{\mathfrak m_3\},
 \qquad f=F.
\]
The source is a regular local ring, so its
puncture has trivial Picard group
\citep[Lemma 52.27.5, Tag 0F2H]{D_StacksRegularPuncture}.
Strict henselization commutes with this quotient,
and therefore $Y_3=X_3\cap V(F)$ is exactly $X_2$.
The nonzero class $\rho^*Q\in J^1(Y_3)$ has no
preimage from $J^1(X_3)=0$.

Finally, for $N=1$ a normal local domain is a
discrete valuation ring: its puncture is the
spectrum of its fraction field, and the punctured
divisor is empty. Both Picard groups vanish.
For $N=2$, every regular punctured divisor is
a zero-dimensional noetherian scheme and hence
a finite disjoint union of spectra of fields.
Its Picard group is zero, so degree-one
surjectivity is automatic. Together with C.3m--n,
these facts give exactly the dimension claims
stated above.
\end{proof}

These are statements over $\mathbb C$, not a
claim that each failure persists over every
fixed positive-characteristic field. For example,
over $k=\overline{\mathbb F}_p$ every point of an
abelian variety is torsion. The actual
abelian-parameter construction in step 1 of C.3n
then shows that ordinary $\operatorname{Pic}_{\mathrm{alg}}$
of any connected regular noetherian $k$-scheme
is torsion: every class lies in the image of
the point group of an abelian variety. Thus
the rational ordinary $J^1$ groups in this
setting vanish. The characteristic-zero
counterexamples above must not be used to
infer the corresponding nonvanishing over
$\overline{\mathbb F}_p$.

The same torsion conclusion holds for ordinary
$\operatorname{CH}^r_{\mathrm{alg}}$ on these
excellent regular punctures in every codimension.
Here we use Chow groups for schemes of finite
type over the universally catenary base $X$
with its codimension dimension function, so
$X$ need not be of finite type over $k$;
this is precisely the setup of
\citet[Situation 42.7.1 and Lemma 42.20.3]{D_StacksChowProper}.
For a generator represented by a cycle family
over a smooth connected curve $T$ and points
$t_0,t_1\in T(k)$, complete $T$ to a smooth
projective curve $C$ and close each horizontal
integral component in $X\times C$. The closure
is still dominant, hence flat, over the regular
curve. Since $[t_1]-[t_0]$ is torsion in
$\operatorname{Pic}^0(C)(\overline{\mathbb F}_p)$,
there are $a>0$ and $u\in k(C)^\times$ with
\[
                 \operatorname{div}_C(u)
                         =a([t_1]-[t_0]).
\]
On each horizontal component, the pulled-back
rational function has divisor equal to $a$
times the difference of its two fibre cycles,
with their Cartier-divisor multiplicities.
The projection of that closure to $X$ is
proper and of finite type, so its pushforward
is rationally equivalent to zero by the cited
lemma. Thus $a$ kills the given fibre difference.
Finite sums are torsion as well. This proves
$\operatorname{CH}^r_{\mathrm{alg}}(X)\otimes\mathbb Q=0$
in this ordinary curve-family convention,
without asserting any relation to the
identity component of a local cycle functor.
